\documentclass[../../main.tex]{subfiles}
\begin{document}
\section{习题}\label{section:ch18-problems}

\begin{enumerate}
\item\label{exercise:ch18-02} 设 $M$ 是一个定向光滑流形，$\omega$ 是 $M$ 上的一个闭 $p$-形式。
\begin{enumerate}[(a)]
\item 证明 $\omega$ 是恰当的当且仅当 $\omega$ 在每个光滑 $p$-闭链上的积分为零。
\item 设 $H_p(M)$ 由有限个光滑 $p$-闭链 $\{c_1,\dots,c_m\}$ 的同调类生成。定义 $\omega$ 关于这组生成元的\emph{周期} $P_i(\omega)=\int_{c_i}\omega$。证明 $\omega$ 恰当当且仅当它的所有周期都为零。
\end{enumerate}

\item\label{exercise:ch18-03} 设 $G$ 是一个群，$[G,G]$ 是其换位子群，$\operatorname{Ab}(G)=G/[G,G]$ 是其 Abel 化。设 $M$ 是一个连通光滑流形，$q\in M$。已知存在一个从 $\pi_1(M,q)$ 到 $H_1(M)$ 的群同态，它把环路 $\gamma$ 的同伦类映为由 $\gamma$ 确定的 $1$-闭链的同调类，并下降为从 $\operatorname{Ab}(\pi_1(M,q))$ 到 $H_1(M)$ 的同构。利用此结果与 de Rham 定理，证明映射 $\Phi:H^1_{\mathrm{dR}}(M)\to\operatorname{Hom}(\pi_1(M,q),\mathbb{R})$ 是一个同构。

\item\label{exercise:ch18-04} 设 $M$ 是一个光滑 $n$ 流形，$S\subseteq M$ 是一个定向紧嵌入 $p$ 维子流形。$S$ 的一个\emph{光滑三角剖分}是 $M$ 中一个光滑 $p$-链 $c=\sum_i\sigma_i$，满足：每个 $\sigma_i:\Delta^p\to S$ 是保定向光滑嵌入；当 $i\neq j$ 时 $\sigma_i(\operatorname{Int}\Delta^p)$ 两两不交；$S=\bigcup_i\sigma_i(\Delta^p)$；且 $\partial c=0$。两个定向紧嵌入 $p$ 维子流形 $S,S'$ 若分别有三角剖分 $c,c'$ 使 $c-c'$ 为边缘，则称它们\emph{同调}。
\begin{enumerate}[(a)]
\item 证明对 $S$ 的任意光滑三角剖分 $c$ 与 $M$ 上任意光滑 $p$-形式 $\omega$，有 $\int_c\omega=\int_S\omega$。
\item 证明：若 $\omega$ 闭且 $S,S'$ 同调，则 $\int_S\omega=\int_{S'}\omega$。
\end{enumerate}

\item\label{exercise:ch18-05} 设 $(M,g)$ 是一个 Riemann $n$ 流形。若 $\omega$ 闭，且只要 $(v_1,\dots,v_p)$ 是某切空间 $T_xM$ 中的标准正交向量就有 $|\omega_x(v_1,\dots,v_p)|\leqslant1$，则称 $\omega$ 是一个\emph{校准}。若存在一个校准 $\omega$ 使 $\eta_S^*\omega$ 是 $S$ 上诱导度量的体积形式，则称定向嵌入 $p$ 维子流形 $S\subseteq M$ 是\emph{被校准的}。设 $S\subseteq M$ 是一个被三角剖分的、被校准的紧子流形。证明 $S$（关于诱导 Riemann 度量）的体积小于或等于任何与 $S$ 同调的子流形的体积。

\item\label{exercise:ch18-06} 设 $D\subseteq\mathbb{R}^3$ 是圆 $(r-2)^2+z^2=1$ 绕 $z$ 轴旋转所得曲面（带诱导 Riemann 度量），$C\subseteq D$ 是内圆 $C=\{(x,y,z):z=0,\ x^2+y^2=1\}$。证明 $C$ 是被校准的，因而在其同调类中是最短曲线。

\item\label{exercise:ch18-07} 对任意光滑流形 $M$，记 $H^p_c(M)$ 为 $M$ 的 $p$ 次紧支集 de Rham 上同调群。
\begin{enumerate}[(a)]
\item 设 $U\subseteq M$ 为开子集，$\eta:U\hookrightarrow M$ 为包含映射。把每个紧支集形式在 $M\setminus U$ 上延拓为零，定义 $\eta_\sharp:\Omega^p_c(U)\to\Omega^p_c(M)$。证明 $\mathrm{d}\circ\eta_\sharp=\eta_\sharp\circ\mathrm{d}$，因而诱导 $\eta^*:H^p_c(U)\to H^p_c(M)$。
\item \textbf{紧支集 Mayer--Vietoris：}设 $M=U\cup V$。证明存在一个连接同态 $\delta^*:H^p_c(M)\to H^{p+1}_c(U\cap V)$ 使相应序列正合。
\end{enumerate}

\item\label{exercise:ch18-08} \textbf{Poincaré 对偶定理：}设 $M$ 是一个定向光滑 $n$ 流形。定义映射 $\mathrm{PD}:\Omega^p(M)\to\Omega^{n-p}_c(M)^*$ 为 $\mathrm{PD}(\omega)(\eta)=\int_M\omega\wedge\eta$。
\begin{enumerate}[(a)]
\item 证明 $\mathrm{PD}$ 下降为线性映射 $\mathrm{PD}:H^p_{\mathrm{dR}}(M)\to H^{n-p}_c(M)^*$。
\item 证明 $\mathrm{PD}$ 对每个 $p$ 都是同构。[提示：模仿 de Rham 定理的证明。]
\end{enumerate}

\item\label{exercise:ch18-09} 设 $M$ 是一个紧光滑 $n$ 流形。
\begin{enumerate}[(a)]
\item 证明 $M$ 的所有 de Rham 群都是有限维的。
\item 证明：若 $M$ 可定向，则对所有 $p$，$\dim H^p_{\mathrm{dR}}(M)=\dim H^{n-p}_{\mathrm{dR}}(M)$。
\end{enumerate}

\item\label{exercise:ch18-10} 设 $M$ 是一个所有 de Rham 群都有限维的光滑 $n$ 流形。$M$ 的\emph{Euler 示性数}是数
\begin{align*}
\chi(M)=\sum_{p=0}^n(-1)^p\dim H^p_{\mathrm{dR}}(M).
\end{align*}
证明 $\chi(M)$ 是 $M$ 的同伦不变量；并证明当 $M$ 紧、可定向、维数为奇数时 $\chi(M)=0$。
\end{enumerate}
\end{document}
